% =====================================================================
%  GVHDNhanXet.tex — Trang NHẬN XÉT CỦA GIÁO VIÊN HƯỚNG DẪN
%
%  Trang standalone trong phần front matter (đánh số La Mã).
%  Bố cục (đã tinh chỉnh để GỌN vừa MỘT trang A4 với 12 dòng viết):
%    - Không đánh số trang (\thispagestyle{empty})
%    - Tiêu đề nhỏ vừa, info block \footnotesize để gọn
%    - Nhãn "Nội dung nhận xét" + 12 dòng kẻ cách nhau 0.9em
%    - Khối ngày/tháng/năm + chữ ký đặt cuối trang
%
%  Trang này được \input từ main.tex; được \clearpage ở đầu và cuối
%  để đảm bảo nằm trên MỘT trang riêng biệt, không tràn.
% =====================================================================

\clearpage
\thispagestyle{empty}

% Lưu ý: \thispagestyle{empty} ở trên đã đủ để ẩn header/footer trên
% TRANG NÀY (GVHD nhận xét). KHÔNG dùng \let\@oddfoot\@empty vì nó
% làm thay đổi page style TOÀN CỤC cho mọi trang phía sau — đã từng
% gây lỗi "mất số trang" từ Chương 1 trở đi. Xem preamble.sty mục
% "QUY ƯỚC ĐÁNH SỐ TRANG" để biết chi tiết.

\begin{center}
    {\normalsize \bfseries NHẬN XÉT CỦA GIÁO VIÊN HƯỚNG DẪN \par}
\end{center}

\vspace{0.4em}

% ---------- Khối thông tin sinh viên + đề tài ----------
% Dùng \footnotesize để khối info gọn, dùng \parbox thu gọn chiều rộng
% của cột "Tên đề tài" để bẻ dòng tự động trong 1 ô tabular.
\noindent\footnotesize
\begin{tabular}{@{}p{4.0cm}@{\hspace{0.3em}}p{11cm}@{}}
    Giảng viên hướng dẫn: & TS. Lê Quang Minh \\
    Sinh viên thực hiện:  & Nguyễn Hoàng Anh \\
    Ngày sinh:            & 27/06/2005 \\
    Lớp:                  & DH23AI01C \\
    Nơi sinh:             & TP Hồ Chí Minh \\
    Tên đề tài:           & PHÁT HIỆN TIN GIẢ TIẾNG VIỆT SỬ DỤNG HỌC MÁY VÀ HỌC SÂU \\
\end{tabular}

\vspace{0.5em}

\noindent Nội dung nhận xét của giảng viên hướng dẫn:

\vspace{0.3em}

% ---------- 12 dòng kẻ cho GVHD viết nhận xét ----------
% Dùng \rule cao 0pt để giữ baseline chính xác + \\ để xuống dòng
% đúng cách. Mỗi dòng cách nhau 1.0em (~1cm với \onehalfspacing).
\noindent\rule{0pt}{0.4pt}\hrulefill\\[1.0em]
\noindent\rule{0pt}{0.4pt}\hrulefill\\[1.0em]
\noindent\rule{0pt}{0.4pt}\hrulefill\\[1.0em]
\noindent\rule{0pt}{0.4pt}\hrulefill\\[1.0em]
\noindent\rule{0pt}{0.4pt}\hrulefill\\[1.0em]
\noindent\rule{0pt}{0.4pt}\hrulefill\\[1.0em]
\noindent\rule{0pt}{0.4pt}\hrulefill\\[1.0em]
\noindent\rule{0pt}{0.4pt}\hrulefill\\[1.0em]
\noindent\rule{0pt}{0.4pt}\hrulefill\\[1.0em]
\noindent\rule{0pt}{0.4pt}\hrulefill\\[1.0em]
\noindent\rule{0pt}{0.4pt}\hrulefill\\[1.0em]
\noindent\rule{0pt}{0.4pt}\hrulefill\par

% ---------- Khối ngày/tháng + chữ ký đặt cách đáy trang, canh phải ----------
% Dùng \vfill để đẩy khối date+signature xuống vùng đáy trang.
\vfill

\begin{flushright}
\begin{minipage}{0.65\linewidth}
Thành phố Hồ Chí Minh, ngày \dots\ tháng \dots\ năm 2026 \\[1.2em]
\hspace*{0.55\linewidth}Người nhận xét
\end{minipage}
\end{flushright}

\clearpage
